\section*{Hoofdstuk \ref{ch:g7:plate-tectonics} --- Platentektoniek: de bewegende puzzel}

\begin{solution}{exo:g7:plate-tectonics:1}
\omterm{def:g7:plate-tectonics:continental-drift}{Continentverschuiving} is het idee dat de \omterm{def:g3:continents-and-oceans:continent}{continenten} langzaam over het
aardoppervlak bewegen, zodat hun plaats in de loop van de tijd veranderd
is. De aanwijzingen van Wegener: het passen van de kustlijnen
(Zuid-Amerika en Afrika); dezelfde fossielen op \omterm{def:g3:continents-and-oceans:continent}{continenten} die nu door
\omterm{def:g2:water-on-earth:ocean}{oceanen} gescheiden zijn; dezelfde \omterm{def:g3:stones-pebbles-sand:rock}{gesteenten} en oude \omterm{def:g4:world-physical-map:range}{gebergten} aan
weerszijden van een \omterm{def:g2:water-on-earth:ocean}{oceaan}; sporen van vroegere klimaten op de verkeerde
plaats (steenkool op Antarctica, gletsjerafzettingen in het hete, droge
Afrika).
\end{solution}

\begin{solution}{exo:g7:plate-tectonics:2}
Divergent: de \omterm{def:g7:plate-tectonics:plate}{platen} bewegen uit elkaar. Convergent: ze bewegen naar elkaar
toe. Transform: ze schuiven zij aan zij langs elkaar.
\end{solution}

\begin{solution}{exo:g7:plate-tectonics:3}
Bijvoorbeeld helemaal rond de Stille \omterm{def:g2:water-on-earth:ocean}{Oceaan} (Japan, Indonesië, de Andes),
en de gordel die van de Middellandse \omterm{def:g2:water-on-earth:ocean}{Zee} tot de Himalaya loopt. Bijna geen
grote \omterm{def:g4:volcanoes-earthquakes-first-look:earthquake}{aardbevingen} in het midden van Afrika of van Australië (ook niet in
Noord-Amerika ten oosten van de Rocky Mountains, of in Brazilië).
\end{solution}

\begin{solution}{exo:g7:plate-tectonics:4}
Een zoetwaterdier kon geen duizenden kilometers zoute \omterm{def:g2:water-on-earth:ocean}{oceaan} oversteken.
Dat het aan beide kanten wordt gevonden, betekent dat de twee plaatsen
vroeger door land met elkaar verbonden waren, in één \omterm{def:g3:continents-and-oceans:continent}{continent}, toen
\emph{Mesosaurus} leefde.
\end{solution}

\begin{solution}{exo:g7:plate-tectonics:5}
Steenkool ontstaat uit de planten van warme, natte, moerassige bossen: het
\omterm{def:g5:weather-and-climate:climate}{klimaat} was warm en vochtig. Zulke bossen kunnen nu niet groeien bij de
\omterm{def:g3:continents-and-oceans:poles}{Zuidpool}, dus moet Antarctica ooit veel dichter bij de \omterm{def:g3:continents-and-oceans:equator}{evenaar} hebben
gelegen: het is verschoven.
\end{solution}

\begin{solution}{exo:g7:plate-tectonics:6}
Midden-Atlantische Rug: divergent. Andes: convergent (een oceanische \omterm{def:g7:plate-tectonics:plate}{plaat}
zakt onder een \omterm{def:g3:continents-and-oceans:continent}{continent}). San Andreasbreuk: transform. Himalaya:
convergent (twee \omterm{def:g3:continents-and-oceans:continent}{continenten} botsen). Oost-Afrikaanse Slenk: divergent (een
\omterm{def:g3:continents-and-oceans:continent}{continent} dat splijt).
\end{solution}

\begin{solution}{exo:g7:plate-tectonics:7}
$2.5 \times 10 = 25$ kilometer per miljoen jaar; in 60 miljoen jaar:
$25 \times 60 = \qty{1500}{km}$.
\end{solution}

\begin{solution}{exo:g7:plate-tectonics:8}
$\qty{7}{m} = \qty{700}{cm}$, en $700 \div 9 \approx 78$: ongeveer
\qty{78}{cm} per jaar, zo'n 30 keer het gemiddelde van \qty{2.5}{cm} per
jaar. De \omterm{def:g7:plate-tectonics:plate}{platen} schuiven aan het oppervlak niet gelijkmatig uit elkaar: de
spanning bouwt jarenlang op en komt vrij in korte uitbarstingen, waarbij
de spleten in één keer meters opengaan. Over lange perioden middelen rustige
jaren en uitbarstingen zich uit tot een paar centimeter per jaar.
\end{solution}

\begin{solution}{exo:g7:plate-tectonics:9}
Hij kon niet zeggen welke kracht de \omterm{def:g3:continents-and-oceans:continent}{continenten} beweegt, en de beweging
die hij zich voorstelde (\omterm{def:g3:continents-and-oceans:continent}{continenten} die door de oceaanbodem ploegen) was
onmogelijk. In de \omterm{def:g7:plate-tectonics:plate-tectonics}{platentektoniek} ploegen de \omterm{def:g3:continents-and-oceans:continent}{continenten} nergens doorheen:
ze worden meegedragen, samen met de oceaanbodem eromheen, door \omterm{def:g7:plate-tectonics:plate}{platen} die
als hele stukken van de buitenste schil van de Aarde bewegen.
\end{solution}

\begin{solution}{exo:g7:plate-tectonics:10}
Het hoogst bij de rug, het diepst bij de trog. Nieuwe \omterm{def:g7:plate-tectonics:plate}{plaat} wordt gemaakt
bij de rug, waar heet \omterm{def:g3:stones-pebbles-sand:rock}{gesteente} uit de \omterm{def:g6:inside-the-earth:layers}{mantel} opstijgt; de oceanische \omterm{def:g7:plate-tectonics:plate}{plaat}
gaat bij de trog terug de \omterm{def:g6:inside-the-earth:layers}{mantel} in, waar ze onder het \omterm{def:g3:continents-and-oceans:continent}{continent} zakt.
\end{solution}

\begin{solution}{exo:g7:plate-tectonics:11}
$14.2 \div 2.2 \approx 6.45$: ongeveer 6 keer sneller. In
kilometer per miljoen jaar zijn de snelheden 142 en 22. Tijden voor
\qty{1000}{km}: $1000 \div 142 \approx 7$ miljoen jaar voor de
Oost-Pacifische Rug; $1000 \div 22 \approx 45$ miljoen jaar voor de
Noord-Atlantische rug.
\end{solution}

\begin{solution}{exo:g7:plate-tectonics:12}
(a) $6472 - 4578 = 1894$ \omterm{def:g4:volcanoes-earthquakes-first-look:earthquake}{aardbevingen}. (b) $4578 \div 7031 \approx 0.651$,
dat is 65\,\%. (c) De lijn op de \omterm{def:g2:drawing-to-map:map}{kaart} geeft de trog aan, waar de
oceanische \omterm{def:g7:plate-tectonics:plate}{plaat} begint te zakken. De \omterm{def:g7:plate-tectonics:plate}{plaat} gaat niet recht naar beneden:
ze zakt schuin onder de andere \omterm{def:g7:plate-tectonics:plate}{plaat}, steeds dieper naarmate ze eronder
verdwijnt. Langs de hele zakkende \omterm{def:g7:plate-tectonics:plate}{plaat} vinden \omterm{def:g4:volcanoes-earthquakes-first-look:earthquake}{aardbevingen} plaats, dus
liggen er veel onder de bovenliggende \omterm{def:g7:plate-tectonics:plate}{plaat}, aan dezelfde kant van de
trog, tientallen of honderden kilometers van de lijn. Een \omterm{def:g7:plate-tectonics:plate-boundary}{plaatgrens} is
een zone, geen dunne lijn.
\end{solution}

\begin{solution}{pb:g7:plate-tectonics:1}
\textbf{1.} \emph{Mesosaurus}, \emph{Cynognathus} en \emph{Glossopteris}
(de band van \emph{Lystrosaurus} verbindt Afrika, India en Antarctica).

\textbf{2.} Het leefde in \omterm{def:g2:water-on-earth:salt-fresh}{zoet water}: het kon geen zoute \omterm{def:g2:water-on-earth:ocean}{oceaan} van
duizenden kilometers breed oversteken.

\textbf{3.} De bolling van het noordoosten van Brazilië past in de holte
van de Golf van Guinee, aan de westkust van Afrika.

\textbf{4.} Een \omterm{def:g7:plate-tectonics:plate-boundary}{divergente grens}, de Midden-Atlantische Rug: de
Zuid-Amerikaanse \omterm{def:g7:plate-tectonics:plate}{plaat} beweegt naar het westen en de Afrikaanse \omterm{def:g7:plate-tectonics:plate}{plaat} naar
het oosten, van elkaar af.

\textbf{5.} De \omterm{def:g1:landscapes-around-us:water}{kusten} liggen aan weerszijden van de \omterm{def:g6:latitude-longitude:longitude}{nulmeridiaan}:
$43.2 + 14.4 = \ang{57.6}$.

\textbf{6.} $57.6 \times 102 = 5875.2$: ongeveer \qty{5880}{km}.

\textbf{7.} De \omterm{def:g6:latitude-longitude:longitude}{meridianen} komen dichter bij elkaar naarmate je van de
\omterm{def:g3:continents-and-oceans:equator}{evenaar} naar een \omterm{def:g3:continents-and-oceans:poles}{pool} gaat. De \omterm{def:g6:latitude-longitude:latitude}{breedtecirkel} van \ang{23} \omterm{def:g6:latitude-longitude:latitude}{zuiderbreedte} is
een kleinere cirkel dan de \omterm{def:g3:continents-and-oceans:equator}{evenaar}, dus is elk van zijn 360 graden korter.

\textbf{8.} Ongeveer de helft van de breedte: $5880 \div 2 = \qty{2940}{km}$.

\textbf{9.} $3.4 \times 10 = 34$ kilometer per miljoen jaar.

\textbf{10.} $5880 \div 34 \approx 173$: ongeveer 170 miljoen jaar.

\textbf{11.} Ja: 170 miljoen jaar ligt tussen 100 en 200 miljoen jaar, en
het komt na het begin van het uiteenvallen van \omterm{def:g7:plate-tectonics:supercontinent}{Pangea}, 225 tot 200 miljoen
jaar geleden. Een eenvoudige berekening op een \omterm{def:g2:drawing-to-map:map}{kaart} geeft de juiste orde
van grootte.

\textbf{12.} Bijvoorbeeld: de snelheid is niet altijd dezelfde geweest; de
\omterm{def:g3:continents-and-oceans:continent}{continenten} werden uitgerekt en dunner voordat ze spleten, dus de huidige
\omterm{def:g1:landscapes-around-us:water}{kusten} zijn niet precies de lijn waarlangs ze uit elkaar braken.
\end{solution}
